% Part: first-order-logic
% Chapter: syntax-and-semantics
% Section: formation-sequences

\documentclass[../../../include/open-logic-section]{subfiles}

\begin{document}

\olfileid[hi]{fol}{syn}{fseq}

\olsection{निर्माण-अनुक्रम}

!!{formula}s को आगमनात्मक परिभाषा द्वारा परिभाषित करना और उसकी पूरक तकनीक,
अर्थात् !!{formula}s के गुणधर्म आगमन द्वारा सिद्ध करना, सुरुचिपूर्ण और दक्ष
दृष्टिकोण है। फिर भी !!{formula}s के निर्माण के अधिक आधारोन्मुख, चरण-दर-चरण
दृष्टिकोण पर विचार करना भी उपयोगी हो सकता है। यहाँ हम \emph{निर्माण-अनुक्रम}
की धारणा से ऐसा करेंगे।
%
यह दिखाने के लिए कि पदों और !!{formula}s को उनकी आगमनात्मक परिभाषाओं का
उल्लेख किए बिना इस ढंग से कैसे प्रस्तुत किया जा सकता है, हम पहले किसी
भाषा~$\Lang L$ से लिए प्रतीकों की मनमानी शृंखला की धारणा प्रस्तुत करते हैं।

\begin{defn}[शृंखलाएँ]
\ollabel{defn:string}
मान लें कि $\Lang L$ कोई प्रथम-क्रम भाषा है। \emph{$\Lang
L$-शृंखला},
~$\Lang L$ के प्रतीकों का कोई परिमित अनुक्रम है। जहाँ भाषा~$\Lang L$ संदर्भ
से स्पष्टतः नियत हो, वहाँ हम $\Lang L$-शृंखला को प्रायः केवल \emph{शृंखला}
कहेंगे।
\end{defn}

\begin{ex}
किसी भी प्रथम-क्रम भाषा $\Lang L$ के लिए सभी $\Lang L$-!!{formula}s,
$\Lang L$-शृंखलाएँ हैं, किंतु इसका विलोम सत्य नहीं है। उदाहरणतः
\[)(\Obj v_0\lif\lexists\] एक $\Lang L$-शृंखला है, पर $\Lang L$-!!{formula}
नहीं है।
\end{ex}

\begin{defn}[पदों के निर्माण-अनुक्रम]
\ollabel{defn:fseq-trm}
$\Lang L$-शृंखलाओं का कोई परिमित अनुक्रम $\tuple{t_0,\dotsc,t_n}$ किसी
पद $t$ का \emph{निर्माण-अनुक्रम} है, यदि $t \ident t_n$ और प्रत्येक $i \leq n$
के लिए या तो $t_i$ !!a{variable} अथवा !!a{constant} है, या~$\Lang
L$
में कोई $k$-स्थानी !!{function}~$f$ है और ऐसे $m_0,\dotsc,m_k < i$ हैं कि
$t_i \ident f(t_{m_0},\dotsc,t_{m_k})$। जहाँ भेद करना आवश्यक हो, वहाँ पदों
के निर्माण-अनुक्रमों को हम \emph{पद-निर्माण-अनुक्रम} कहेंगे।
\end{defn}

\begin{ex}
अनुक्रम
\[
    \tuple{\Obj c_0, \Obj v_0, \Atom{\Obj f^2_0}{\Obj c_0, \Obj v_0}, \Atom{\Obj f^1_0}{\Atom{\Obj f^2_0}{\Obj c_0, \Obj v_0}}}
\]
पद $\Atom{\Obj f^1_0}{\Atom{\Obj
f^2_0}{\Obj c_0, \Obj v_0}}$ का एक
निर्माण-अनुक्रम है; इसी प्रकार यह भी है:
\[
    \tuple{\Obj v_0, \Obj c_0, \Atom{\Obj f^2_0}{\Obj c_0, \Obj v_0}, \Atom{\Obj f^1_0}{\Atom{\Obj f^2_0}{\Obj c_0, \Obj v_0}}}.
\]
\end{ex}

\begin{defn}[सूत्रों के निर्माण-अनुक्रम]
\ollabel{defn:fseq-frm}
$\Lang L$-शृंखलाओं का कोई परिमित अनुक्रम $\tuple{!A_0,\dotsc,!A_n}$,
~$!A$ का \emph{निर्माण-अनुक्रम} है, यदि $!A \ident !A_n$ और प्रत्येक $i \leq n$
के लिए या तो $!A_i$ कोई परमाण्विक !!{formula} है, अथवा ऐसे $j,k < i$
और कोई !!a{variable}~$x$ हैं कि निम्न में से कोई एक लागू हो:
\begin{enumerate}
    \tagitem{prvNot}{$!A_i \ident \lnot !A_j$.}{}%
    \tagitem{prvAnd}{$!A_i \ident (!A_j \land !A_k)$.}{}%
    \tagitem{prvOr}{$!A_i \ident (!A_j \lor !A_k)$.}{}%
    \tagitem{prvIf}{$!A_i \ident (!A_j \lif !A_k)$.}{}%
    \tagitem{prvIff}{$!A_i \ident (!A_j \liff !A_k)$.}{}%
    \tagitem{prvAll}{$!A_i \ident \lforall[x][!A_j]$.}{}%
    \tagitem{prvEx}{$!A_i \ident \lexists[x][!A_j]$.}{}%
\end{enumerate}
जहाँ भेद करना आवश्यक हो, वहाँ सूत्रों के निर्माण-अनुक्रमों को हम
\emph{सूत्र-निर्माण-अनुक्रम} कहेंगे।
\end{defn}

\begin{ex}
\[
\tuple{
    \Atom{\Obj A^1_0}{\Obj v_0},
    \Atom{\Obj A^1_1}{\Obj c_1},
    (\Atom{\Obj A^1_1}{\Obj c_1} \land \Atom{\Obj A^1_0}{\Obj v_0}),
    \lexists[\Obj v_0][(\Atom{\Obj A^1_1}{\Obj c_1} \land \Atom{\Obj A^1_0}{\Obj v_0})]
}
\]
$\lexists[\Obj v_0][(\Atom{\Obj A^1_1}{\Obj
c_1} \land \Atom{\Obj A^1_0}{\Obj v_0})]$ का एक निर्माण-अनुक्रम है; इसी प्रकार यह भी है:
\begin{multline*}
\tuple{
    \Atom{\Obj A^1_0}{\Obj v_0},
    \Atom{\Obj A^1_1}{\Obj c_1},
    (\Atom{\Obj A^1_1}{\Obj c_1} \land \Atom{\Obj A^1_0}{\Obj v_0}),
    \Atom{\Obj A^1_1}{\Obj c_1},\\
    \lforall[\Obj v_1][\Atom{\Obj A^1_0}{\Obj v_0}],
    \lexists[\Obj v_0][(\Atom{\Obj A^1_1}{\Obj c_1} \land \Atom{\Obj A^1_0}{\Obj v_0})]
}.
\end{multline*}
%
दूसरे उदाहरण से स्पष्ट है कि निर्माण-अनुक्रम में ``अनुपयोगी सामग्री'' आ सकती
है: ऐसे !!{formula}s जो अनावश्यक हों अथवा निर्माण में कोई योगदान न दें।
\end{ex}

\begin{prop}\ollabel{prop:formed}
प्रत्येक !!{formula}~$!A$, जो~$\Frm[L]$ में है, का कोई निर्माण-अनुक्रम होता है।
\end{prop}

\begin{proof}
मान लें कि $!A$ परमाण्विक है। तब अनुक्रम $\tuple{!A}$,~$!A$ का
निर्माण-अनुक्रम है।
%
अब मान लें कि $!B$ और~$!C$ के निर्माण-अनुक्रम क्रमशः
$\tuple{!B_0,\dotsc,!B_n}$ और $\tuple{!C_0,\dotsc,!C_m}$ हैं।
%
\begin{enumerate}
  \tagitem{prvNot}{यदि $!A \ident \lnot !B$, तो
      $\tuple{!B_0,\dotsc,!B_n,\lnot !B_n}$,~$!A$ का निर्माण-अनुक्रम है।}{}
  \tagitem{prvAnd}{यदि $!A \ident (!B \land !C)$, तो
      $\tuple{!B_0,\dotsc,!B_n,!C_0,\dotsc,!C_m,(!B_n \land !C_m)}$,
      ~ $!A$ का निर्माण-अनुक्रम है।}{}
  \tagitem{prvOr}{यदि $!A \ident (!B \lor !C)$, तो
      $\tuple{!B_0,\dotsc,!B_n,!C_0,\dotsc,!C_m,(!B_n \lor !C_m)}$,
      ~ $!A$ का निर्माण-अनुक्रम है।}{}
  \tagitem{prvIf}{यदि $!A \ident (!B \lif !C)$, तो
      $\tuple{!B_0,\dotsc,!B_n,!C_0,\dotsc,!C_m,(!B_n \lif !C_m)}$,
      ~ $!A$ का निर्माण-अनुक्रम है।}{}
  \tagitem{prvIff}{यदि $!A \ident (!B \liff !C)$, तो
      $\tuple{!B_0,\dotsc,!B_n,!C_0,\dotsc,!C_m,(!B_n \liff !C_m)}$,
      ~ $!A$ का निर्माण-अनुक्रम है।}{}
  \tagitem{prvAll}{यदि $!A \ident \lforall[x][!B]$, तो
      $\tuple{!B_0,\dotsc,!B_n,\lforall[x][!B_n]}$,~$!A$ का
      निर्माण-अनुक्रम है।}{}
  \tagitem{prvEx}{यदि $!A \ident \lexists[x][!B]$, तो
      $\tuple{!B_0,\dotsc,!B_n,\lexists[x][!B_n]}$,~$!A$ का
      निर्माण-अनुक्रम है।}{}
\end{enumerate}
!!{formula}s पर आगमन के सिद्धांत से प्रत्येक !!{formula} का कोई
निर्माण-अनुक्रम है।
\end{proof}

हम इसका विलोम भी सिद्ध कर सकते हैं। यह महत्त्वपूर्ण है, क्योंकि इससे पता
चलता है कि सूत्रों को परिभाषित करने के हमारे दोनों ढंग तुल्य हैं: वे एक ही
परिणाम देते हैं। इसका यह अर्थ भी है कि निर्माण-अनुक्रमों की लंबाई पर साधारण
आगमन का उपयोग करके हम सूत्रों के बारे में प्रमेय सिद्ध कर सकते हैं।

\begin{lem}
\ollabel{lem:fseq-init}
मान लें कि $\tuple{!A_0,\dotsc,!A_n}$,~$!A_n$ का निर्माण-अनुक्रम है और $k \leq n$।
तब $\tuple{!A_0,\dotsc,!A_k}$,~$!A_k$ का निर्माण-अनुक्रम है।
\end{lem}

\begin{proof}
अभ्यास।
\end{proof}

\begin{prob}
\olref[fol][syn][fseq]{lem:fseq-init} सिद्ध कीजिए।
\end{prob}

\begin{thm}
\ollabel{thm:fseq-frm-equiv}
$\Frm[L]$ उन सभी $\Lang L$-शृंखलाओं~$!A$ का समुच्चय है जिनके लिए~$!A$ का
कोई सूत्र-निर्माण-अनुक्रम अस्तित्व रखता है।
\end{thm}

\begin{proof}
मान लें कि $F$, भाषा~$\Lang L$ की उन सभी प्रतीक-शृंखलाओं का समुच्चय है जिनका
कोई निर्माण-अनुक्रम है। \olref[fol][syn][fseq]{prop:formed} में हम देख चुके
हैं कि $\Frm[L] \subseteq F$, अतः अब हम विलोम सिद्ध करते हैं।

मान लें कि $!A$ का निर्माण-अनुक्रम $\tuple{!A_0,\dotsc,!A_n}$ है। हम सिद्ध
करते हैं कि $!A \in \Frm[L]$; प्रमाण~$n$ पर प्रबल आगमन द्वारा है। हमारी आगमनात्मक
परिकल्पना है कि लंबाई $m < n$ का निर्माण-अनुक्रम रखने वाली प्रत्येक
प्रतीक-शृंखला $\Frm[L]$ में है। निर्माण-अनुक्रम की परिभाषा से या तो $!A \ident !A_n$
परमाण्विक है, या ऐसे $j,k < n$ होने चाहिए जिनके लिए निम्न में
से कोई एक स्थिति लागू हो:
\begin{enumerate}
    \tagitem{prvNot}{$!A \ident \lnot !A_j$.}{}%
    \tagitem{prvAnd}{$!A \ident (!A_j \land !A_k)$.}{}%
    \tagitem{prvOr}{$!A \ident (!A_j \lor !A_k)$.}{}%
    \tagitem{prvIf}{$!A \ident (!A_j \lif !A_k)$.}{}%
    \tagitem{prvIff}{$!A \ident (!A_j \liff !A_k)$.}{}%
    \tagitem{prvAll}{$!A \ident \lforall[x][!A_j]$.}{}%
    \tagitem{prvEx}{$!A \ident \lexists[x][!A_j]$.}{}%
\end{enumerate}
अब हम मामलों के अनुसार तर्क करते हैं। यदि $!A$ परमाण्विक है, तो $!A_n \in \Frm[L_0]$।
इसके बजाय मान लें कि $!A \equiv (!A_j \land !A_k)$।
\olref[fol][syn][fseq]{lem:fseq-init} से $\tuple{!A_0,\dotsc,!A_j}$ और
$\tuple{!A_0,\dotsc,!A_k}$ क्रमशः $!A_j$ और~$!A_k$ के निर्माण-अनुक्रम हैं।
चूँकि ये~$!A$ के निर्माण-अनुक्रम के उचित आरंभिक उपानुक्रम हैं, दोनों की लंबाई
~$n$ से कम है। अतः आगमनात्मक परिकल्पना से $!A_j$ और~$!A_k$,
~$\Frm[L_0]$ में हैं और !!a{formula} की परिभाषा से $(!A_j \land !A_k)$ भी
उसी में है। अन्य मामले समानांतर तर्क से प्राप्त होते हैं।
\end{proof}

पदों के निर्माण-अनुक्रमों के गुण !!{formula}s के निर्माण-अनुक्रमों के गुणों
के समान हैं।

\begin{prop}
\ollabel{prop:fseq-trm-equiv}
$\Trm[L]$ उन सभी $\Lang L$-शृंखलाओं $t$ का समुच्चय है जिनके लिए~$t$ का
कोई पद-निर्माण-अनुक्रम अस्तित्व रखता है।
\end{prop}

\begin{proof}
अभ्यास।
\end{proof}

\begin{prob}
\olref[fol][syn][fseq]{prop:fseq-trm-equiv} सिद्ध कीजिए। संकेत:
\olref[fol][syn][fseq]{thm:fseq-frm-equiv} के प्रमाण में प्रयुक्त रणनीति के
समान रणनीति अपनाइए।
\end{prob}

निर्माण-अनुक्रमों में दो प्रकार की ``अनुपयोगी सामग्री'' आ सकती है: दोहराए गए
अवयव, और ऐसे अवयव जो सूत्र या पद के निर्माण से असंबद्ध हों। न्यूनतम
निर्माण-अनुक्रमों पर विचार करके हम दोनों को हटा सकते हैं।

\begin{defn}[न्यूनतम निर्माण-अनुक्रम]
\ollabel{defn:minimal-fseq}
कोई निर्माण-अनुक्रम $\tuple{!A_0, \ldots, !A_n}$ किसी सूत्र~$!A$ का है। वह
~$!A$ का \emph{न्यूनतम निर्माण-अनुक्रम} है, यदि प्रत्येक अन्य निर्माण-अनुक्रम~$s$,
जो~$!A$ का हो, के लिए~$s$ की लंबाई~$n+1$ से अधिक या उसके बराबर हो।

इसी प्रकार, कोई निर्माण-अनुक्रम $\tuple{t_0, \ldots, t_n}$ किसी पद~$t$ का है।
वह~$t$ का \emph{न्यूनतम निर्माण-अनुक्रम} है, यदि प्रत्येक अन्य निर्माण-अनुक्रम~$s$,
जो~$t$ का हो, के लिए~$s$ की लंबाई~$n+1$ से अधिक या उसके बराबर हो।
\end{defn}

ध्यान दें कि किसी सूत्र या पद के एक से अधिक न्यूनतम निर्माण-अनुक्रम हो सकते
हैं, किंतु उनमें ठीक वही शृंखलाएँ होंगी।

\begin{prop}
\ollabel{prop:subformula-equivs}
निम्न कथन तुल्य हैं:
\begin{enumerate}
    \item $!B$,~$!A$ का उप-!!{formula} है।
    \item $!B$,~$!A$ के प्रत्येक निर्माण-अनुक्रम में आता है।
    \item $!B$,~$!A$ के किसी न्यूनतम निर्माण-अनुक्रम में आता है।
\end{enumerate}
\end{prop}

\begin{proof}
अभ्यास।
\end{proof}

\begin{prob}
\olref[fol][syn][fseq]{prop:subformula-equivs} सिद्ध कीजिए।
\end{prob}

\begin{history}
निर्माण-अनुक्रमों को रेमंड स्मलियन ने अपनी पाठ्यपुस्तक \emph{First-Order
Logic} \citep{Smullyan1968} में प्रस्तुत किया था। निर्माण-अनुक्रमों के
अतिरिक्त गुणधर्मों के प्रमाण \citet{Zuckerman1973} में मिलते हैं।
\end{history}

\end{document}
